\documentclass[10pt,a4paper]{amsart}
\usepackage[margin=19mm]{geometry}
\usepackage{amsmath,amssymb,mathtools}
\usepackage[scr=rsfs,cal=euler]{mathalfa}
\usepackage{xfrac}
\usepackage[unicode,hidelinks,bookmarks=false]{hyperref}
\newcommand{\ie}{i.e.\ }
\newcommand{\cf}{cf.\ }
\newcommand{\resp}{respectively\ }
\newcommand{\Subsection}{\subsection}
\providecommand{\nobreakdash}{\nobreak\hbox{-}}
\setlength{\emergencystretch}{2em}
\multlinegap0pt
\usepackage{bm}
\usepackage{bbm}
\usepackage{dsfont}
\usepackage{xspace}
\usepackage{cancel}
\usepackage{mathtools}
\usepackage{amsthm}
\makeatletter
\@ifundefined{theorem}{\newtheorem{theorem}{Theorem}}{}
\@ifundefined{lemma}{\newtheorem{lemma}[theorem]{Lemma}}{}
\makeatother
\title[Generalizations of interval and proper interval graphs for s]{Generalizations of interval and proper interval graphs for simplicial complexes}
\author{Fahimeh Khosh-Ahang Ghasr}
\date{Source: arXiv:2510.16722v1 (CC BY 4.0)}
\begin{document}
\maketitle

\noindent\emph{Source and licence.} Generalizations of interval and proper interval graphs for simplicial complexes. Version arXiv:2510.16722v1 (CC BY 4.0), distributed under the Creative Commons Attribution 4.0 International licence, as recorded in the arXiv licence metadata for this exact version and shown on the abstract page https://arxiv.org/abs/2510.16722v1. Full rights notice: licenses/arxiv-2510.16722-CC-BY-4.0.txt.\par\smallskip

\noindent\textit{Editorial scope.} Counted unit: the Theorem whose source label is \texttt{1.4} in the article (source0.tex lines 226--234, source label \texttt{1.4}), delivered together with the article's own complete original proof of it (source0.tex lines 235--263, one \texttt{proof} environment of 29 lines). No mathematics is written, repaired or re-derived: statement and proof are the article's own text line for line. Required context retained uncounted: global (lines 152--154). Every definition, display and label of those lines is retained unchanged. Omitted from this package and not redistributed: 1--151, 155--225, 264--479. Nothing inside a retained excerpt is omitted or altered: every retained body is delivered exactly as the article's lines are written, so no omission receipt of any kind accompanies this package. Citations: the retained excerpts carry no citation command at all, so no bibliography entry of the article is transcribed and no reference is redistributed. Reference adaptation: none was needed and none was made; every reference inside the delivered unit resolves inside it, so no unresolved reference or citation is produced. Printed slips kept literal: no printed word, symbol, hypothesis or number of the counted statement or of its proof was altered. Layout and class: the article's own class is \texttt{amsart} with packages including amsbsy, amsfonts, amsmath, amssymb, amscd, amsthm, mathrsfs, verbatim, calc, graphicx, epsfig, color, enumitem, geometry; the wrapper is the lane's pinned \texttt{amsart} editorial wrapper at 10pt on A4, which restates the theorem-like environments the retained text needs and adds the article's own macro definitions that the retained text needs (none), with extra wrapper packages (bm, bbm, dsfont, xspace, cancel, mathtools). The delivered numbering is the unit's own. Assets: no retained span contains a figure environment, a graphics-inclusion command, a PDF-page inclusion, a table or a TikZ picture; no image file is copied into this package, the package declares no asset dependency and no image file is redistributed with it. Licence: version arXiv:2510.16722v1 of this article is distributed under the Creative Commons Attribution 4.0 International licence, as recorded in the arXiv licence metadata for this exact version and shown on the abstract page https://arxiv.org/abs/2510.16722v1. A full rights notice is delivered at licenses/arxiv-2510.16722-CC-BY-4.0.txt. Characters: every retained excerpt is pure ASCII, so the delivered engine, XeLaTeX with its default Latin Modern text font, reports no missing character.
\begin{lemma}\label{global}
A $d$-complex $\Delta$ is under closed if and only if there is a labelling $[n]$ on $V(\Delta)$ such that  for each facet $F: i_1\dots i_{d+1}$ and each $j$ with $i_1\leq j\leq i_{d+1}, j\notin F$, we have $\{i_1, \dots , i_d, j\}\in \mathcal{F}(\Delta)$.
\end{lemma}

\begin{theorem}\label{1.4}
\begin{enumerate}
\item Suppose that $G$ is a $d$-strong interval (resp. $d$-strong unit interval, $d$-strong proper interval) graph with no isolated vertex. Then it is $k$-strong interval (resp. $d$-strong unit interval, $d$-strong proper interval) for all $k\geq d$ with the same interval representation.

\item Suppose that $G$ is a $d$-under closed graph. Then it is $k$-under closed for all $k\geq d$ with the same labelling on $V(G)$.

\item Suppose $d>1$ and $G$ is a $d$-unit interval graph. Then it is $k$-unit interval for all $k\geq d$ with the same labelling on $V(G)$.
\end{enumerate}
\end{theorem}

\begin{proof}
It is enough to prove the results for $k=d+1$ when $d+2\geq |V(G)|$.

(1) Suppose $G$ is a $d$-strong interval graph on $[n]$ for which $\{I_i\}_{i\in [n]}$ is an interval representation. Assume that $1\leq i_1< \dots < i_{d+2}\leq n$. We should prove that $G[i_1, \dots , i_{d+2}]$ is connected if and only if $\bigcup_{1\leq \ell \leq d+2}I_{i_\ell}$ is an interval.

$(\Longrightarrow )$ If $G[i_1, \dots , i_{d+2}]$ is connected, then it has at least two distinct noncut vertices, say $i_k$ and $i_{k'}$  (we point out that each vertex in cycles is noncut, also if $G$ has no cycle, it should have at least two leaves which are certainly noncut vertices). Hence $G[i_1, \dots ,\widehat{i_k}, \dots, i_{d+2}]$ and $G[i_1, \dots ,\widehat{i_{k'}}, \dots, i_{d+2}]$ are connected. Since $d\in \mathbb{N}$ and $G$ is $d$-strong interval,  $\bigcup_{1\leq \ell \leq d+2, \ell\neq k}I_{i_\ell}$ and $\bigcup_{1\leq \ell \leq d+2, \ell\neq k'}I_{i_\ell}$ are some intervals with nonempty intersection. Hence, $\bigcup_{1\leq \ell \leq d+2}I_{i_\ell}=(\bigcup_{1\leq \ell \leq d+2, \ell\neq k}I_{i_\ell}) \cup (\bigcup_{1\leq \ell \leq d+2, \ell\neq k'}I_{i_\ell})$ is also an interval.

$(\Longleftarrow )$ If $\bigcup_{1\leq \ell \leq d+2}I_{i_\ell}$ is an interval, then suppose $\mathrm{left} \bigcup_{1\leq \ell \leq d+2}I_{i_\ell}=\mathrm{left}I_{i_k}$ and  $\mathrm{right} \bigcup_{1\leq \ell \leq d+2}I_{i_\ell}=\mathrm{right}I_{i_{k'}}$ for some $1\leq k,k'\leq d+2$. Now if $k=k'$, then since $d+2\geq 3$, $I_{i_k}$ contains $I_{i_s}$ and $I_{i_{s'}}$ for distinct integers $s$ and $s'$ with $s,s'\in \{1, \dots, \widehat{k}, \dots , d+2\}$. Thus $\bigcup_{1\leq \ell \leq d+2, \ell\neq s}I_{i_\ell}$ and $\bigcup_{1\leq \ell \leq d+2, \ell\neq s'}I_{i_\ell}$ are intervals. Also if $k\neq k'$, then either $\bigcup_{1\leq \ell \leq d+2, \ell\neq k}I_{i_\ell}$ is an interval or $I_{i_k}$ contains $I_{i_s}$ for some $1\leq s \leq d+2$ with $s\neq k$. Similarly either $\bigcup_{1\leq \ell \leq d+2, \ell\neq k'}I_{i_\ell}$ is an interval or $I_{i_{k'}}$ contains $I_{i_{s'}}$ for some $1\leq s' \leq d+2$ with $s'\neq k'$. Hence in each case $\bigcup_{1\leq \ell \leq d+2, \ell\neq t}I_{i_\ell}$ and $\bigcup_{1\leq \ell \leq d+2, \ell\neq t'}I_{i_\ell}$ are intervals for some distinct integers $t$ and $t'$ with $1\leq t,t'\leq d+2$. Therefore since  $G$ is $d$-strong interval, the graphs $G[i_1, \dots , \widehat{i_t}, \dots , i_{d+2}]$ and $G[i_1, \dots , \widehat{i_{t'}}, \dots , i_{d+2}]$ are connected. Hence $d+2\geq 3$ implies $G[i_1, \dots, i_{d+2}]$  is connected as desired.

(2) Suppose $G$ is a $d$-under closed graph with labelling $[n]$ on $V(G)$. We use the equivalent definition for under closed $d$ complexes presented in Lemma \ref{global}. Consider a connected subgraph of $G$ with vertices $i_1<\dots < i_{d+2}$ and an integer $j$ with $i_1\leq j\leq i_{d+2}$ and $j\notin \{i_1, \dots , i_{d+2}\}$.  Then for connected graph $G[i_1, \dots , i_{d+2}]$ the following cases occur:

$\bullet$ There exists a noncut vertex, say $i_k$, for some $k$ with $1< k <d+2$.
Therefore $G[i_1, \dots , \widehat{i_k}, \dots , i_{d+2}]$ is connected. Since $G$ is $d$-global interval, $G[i_1, \dots , i_{d+1}]$ and\linebreak $G[i_1, \dots , \widehat{i_k}, \dots , i_{d+1}, j]$ are connected. So $G[i_1, \dots , i_{d+1}, j]$ is also connected.

$\bullet$ $i_k$ is a cut vertex for all $k$ with $1<k < d+2$. Hence $G[i_1, \dots , i_{d+1}]$ and $G[i_2, \dots , i_{d+2}]$ are connected. If $i_1<j<i_{d+1}$,  since $G$ is $d$-global interval, $G[i_1, \dots , i_d, j]$ is connected. Hence  $G[i_1, \dots , i_{d+1}, j]$ is also connected. Else $i_2<j<i_{d+2}$. Therefore $G[i_2, \dots , i_{d+1}, j]$ and so $G[i_1, \dots , i_{d+1},j]$ is connected.

(3) Suppose $G$ is a $d$-unit interval graph. Assume $1 \leq i_1 < \dots < i_{d+2} \leq n$ and $i_1 \leq j_1 < \dots < j_{d+2} \leq i_{d+2}$ such that $G[i_1, \dots, i_{d+2}]$ is connected. We prove that $G[j_1, \dots, j_{d+2}]$ is connected. Since $G$ is $d$-unit interval, for any facet ${i_1, \dots, i_{d+1}} \in \mathcal{F}(\Delta_d(G))$, every ${j_1', \dots, j_{d+1}'}$ with $i_1 \leq j_1' < \dots < j_{d+1}' \leq i_{d+1}$ is a facet. We consider two cases for $G[i_1, \dots, i_{d+2}]$.

$\bullet$ There exists a noncut vertex $i_k$, for some $1 < k < d+2$. Thus, $G[i_1, \dots, \widehat{i_k}, \dots, i_{d+2}]$ is connected. Since $G$ is $d$-unit interval and $i_1 \leq j_1 < \dots < j_{d+2} \leq i_{d+2}$, any $d+1$ vertices $\{j_{s_1}, \dots, j_{s_{d+1}}\} \subseteq \{j_1, \dots, j_{d+2}\}$ with $j_{s_1} < \dots < j_{s_{d+1}}$ form a facet of $\Delta_d(G)$. For $d \geq 1$, take any two subsets $\{j_1, \dots, \widehat{j_\ell}, \dots, j_{d+2}\}$ and $\{j_1, \dots, \widehat{j_{\ell'}}, \dots, j_{d+2}\}$ for $\ell \neq \ell'$. These are connected, share $d \geq 1$ vertices, and their union is $G[j_1, \dots, j_{d+2}]$, which is connected.

$\bullet$ Every $i_k$ for $1 < k < d+2$ is a cut vertex. Hence, $G[i_1, \dots, i_{d+2}]$ is a path graph with endpoints $i_1, i_{d+2}$, so $G[i_1, \dots, i_{d+1}]$ and $G[i_2, \dots, i_{d+2}]$ are connected. If $j_{d+2} \leq i_{d+1}$, then $i_1 \leq j_1 < \dots < j_{d+2} \leq i_{d+1}$. Since $G[i_1, \dots, i_{d+1}]$ is connected, $G[j_1, \dots, j_{d+1}]$  and $G[j_2, \dots, j_{d+2}]$ are connected by the $d$-unit interval property.  Their union $G[j_1, \dots, j_{d+2}]$ is connected, as they share $j_2, \dots, j_{d+1}$ and $d \geq 1$.

If $i_{d+1} < j_{d+2} \leq i_{d+2}$, consider $G[i_2, \dots, i_{d+1}, j_{d+2}]$. Since $i_2 \leq i_2 < \dots < i_{d+1} < j_{d+2} \leq i_{d+2}$ and $G[i_2, \dots, i_{d+2}]$ is connected, $G[i_2, \dots, i_{d+1}, j_{d+2}]$ is connected. If $j_1 \geq i_2$, then $i_2 \leq j_1 < \dots < j_{d+2} \leq i_{d+2}$, so $G[j_1, \dots, j_{d+2}]$ is connected. If $i_1 \leq j_1 < i_2$, consider $G[j_1, i_2, \dots, i_{d+1}]$. Since $i_1 \leq j_1 < i_2 < \dots < i_{d+1} \leq i_{d+1}$ and $G[i_1, \dots, i_{d+1}]$ is connected, $G[j_1, i_2, \dots, i_{d+1}]$ is connected. Thus, $G[j_1, i_2, \dots, i_{d+1}, j_{d+2}]$ is connected, as it is the union of $G[j_1, i_2, \dots, i_{d+1}]$ and $G[i_2, \dots, i_{d+1}, j_{d+2}]$, sharing $i_2, \dots, i_{d+1}$.

Now if  $i_k$ is a noncut vertex of $G[j_1, i_2 , \dots , i_{d+1}, j_{d+2}]$ for some $2\leq k \leq d+1$, then similar to the aforementioned case the result holds. Suppose $G[j_1, i_2 , \dots , i_{d+1}, j_{d+2}]$ is a path graph with end vertices $j_1$ and $j_{d+2}$. If $i_2\leq j_2<\dots <j_{d+1} \leq i_{d+1}$, then since $G$ is $d$-unit interval and $G[j_1, i_2 , \dots , i_{d+1}, j_{d+2}]$ is a path graph, we should have $G[j_1, \dots , j_{d+1}]$ and $G[j_2, \dots , j_{d+2}]$ are connected. Their union $G[j_1, \dots, j_{d+2}]$ is connected, as they share $j_2, \dots, j_{d+1}$ and $d \geq 1$. So either $j_2<i_2$ or $i_{d+1}<j_{d+1}$.

Assume that $i_{d+1}<j_{d+1}$ and $i_k$ is the only neighbour of $j_{d+2}$ in $G[j_1, i_2 , \dots , i_{d+1}, j_{d+2}]$. Now $G[i_2, \dots , i_{d+1}, j_{d+2}]$ is connected, $i_2\leq i_2< \dots <\widehat{i_k} < \dots < i_{d+1}<j_{d+1}<j_{d+2}\leq j_{d+2}$ and $d$-unit interval property of $G$ implies $G[ i_2, \dots , \widehat{i_k}, \dots , i_{d+1},j_{d+1},j_{d+2}]$ is connected.  Therefore  $j_{d+1}j_{d+2}\in E(G)$ and $G[ i_2, \dots , \widehat{i_k}, \dots , i_{d+1},j_{d+1}]$ is connected. On the other hand  $G[j_1, i_2, \dots , \widehat{i_k}, \dots , i_{d+1}]$ is connected. So $G[j_1, i_2, \dots , \widehat{i_k}, \dots , i_{d+1},j_{d+1}]$ is  connected. Thus  we have  $G[j_1, \dots , j_{d+1}]$ is connected, because $G$ is $d$-unit interval. This shows $G[j_1, \dots , j_{d+2}]$ is connected as required. The case $j_2<i_2$ is
symmetric.
\end{proof}

\end{document}
